\chapter{病例对照研究的设计}
\chapterauthors{Thomas B. Newman、Warren S. Browner}
\begingroup\setlength{\tabcolsep}{5pt}
\sourcepage{242}
第8章介绍了队列研究：随访一组参与者，估计暴露或未暴露于各种危险因素者的结局发生。相比之下，\term{病例对照研究}从一组已经发生结局者（\term{病例}，cases）和一组没有该结局者（\term{对照}，controls）入手，向后查看，测量两组既往暴露，寻找可解释为何病例出现结局而对照没有的差异。

例如，可以组建眼黑色素瘤（ocular melanoma）病例组与健康对照样本，收集两组既往电弧焊（arc welding）暴露资料，估计其对该病风险的影响（1）。病例对照设计费用相对较低，特别适用于罕见病和疾病暴发研究。

本章还介绍数种变体。\term{巢式病例对照研究}（nested case-control study）将明确队列内出现的新发病例，与从队列其他成员抽取的对照比较。\term{巢式病例队列研究}（nested case-cohort study）以整个队列（包括病例）的随机样本作对照，常可用于数种不同结局。\term{发病密度病例对照设计}（incidence-density case-control design）可考虑危险因素随时间的变化及失访。\term{病例交叉研究}（case-crossover studies）使病例充当自身对照。最后讨论如何在本章和第8章的观察性设计中选择。

\sourcepage{243}
\section{病例对照研究的基础}
多数疾病相对少见，要发现眼黑色素瘤等罕见病的危险因素，一般人群的队列或横断面研究需纳入数千人，费用很高。病例系列有时可借助一般人群危险因素流行程度的既有知识，找到明显危险因素（如艾滋病与注射毒品），但多数因素仍须组建参照组，比较病例与对照的暴露。

按第8章描述的大多数方面，病例对照属于回顾性研究：找出患病与未患病两组，再回头测量疾病发生前的预测因素（图9.1）。不过，研究开始时未必所有结局都已发生。例如，疾病仍在暴发时，流行病学团队即可前往调查，界定病例定义，并随新病例发生前瞻性招募病例与对照，直到找出病因（例9.1）。
\sourcepage{244}
\usetikzlibrary{shapes.geometric}
\providecommand{\chNineCaseControl}{\begin{tikzpicture}[font=\small,>=Stealth]
\foreach \yy/\type/\pop in {0/病例/已发生结局的可接触人群,-4.2/对照/未发生结局的可接触人群}{\begin{scope}[yshift=\yy cm]
\draw[dashed,dcrred] (3,0) ellipse (2 and 1.55);
\node[text=dcrred] at (2.5,1.9) {\pop};
\node[draw,text width=31mm,align=center] (s) at (3,0) {\textcolor{dcrgold}{\type}\\符合选择标准的样本};
\node[draw,text width=29mm,align=center,text=dcrpurple] (m) at (-3,0) {测量预测变量\\（常为既往值）};
\draw[thick,->,dcrteal] (s)--node[above]{回顾}(m);
\end{scope}}
\node[align=center,font=\footnotesize] at (3,-5.4) {对照通常更多，\quad 可与病例匹配};
\node[align=center,font=\footnotesize] at (-3,-5.4) {与病例采用相同测量方式};
\end{tikzpicture}}
\providecommand{\chNineSelection}{\begin{tikzpicture}[font=\small,>=Stealth]
\node at (0,3.1) {疾病新发病例};
\fill[dcrblue!70] (-.35,2.7)--(.35,2.7)--(.35,1.5)--(.18,1.5)--(.18,.4)--(.1,.4)--(.1,-.7)--(.04,-.7)--(.04,-1.8)--(-.04,-1.8)--(-.04,-.7)--(-.1,-.7)--(-.1,.4)--(-.18,.4)--(-.18,1.5)--(-.35,1.5)--cycle;
\foreach \yy/\txt in {1.8/未就医,.7/在别处就诊,-.4/就诊但误诊,-1.5/诊断前死亡或缓解}{\draw[->,thick] (.15,\yy)--(2.1,\yy) node[right]{\txt};}
\draw[->] (0,-1.8)--(0,-2.4);\node at (0,-2.7) {可供病例对照研究的病例};
\end{tikzpicture}}
\providecommand{\chNineNested}{\begin{tikzpicture}[font=\footnotesize,>=Stealth]
\draw[dashed,dcrred] (0,0) ellipse (2.3 and 2.4);\node[text=dcrred] at (0,2.7) {人群};
\node[draw,text width=36mm,align=center,minimum height=28mm] (c) at (0,.5) {既已界定的队列\\预测因素可从保存的\\标本、影像等测量};
\node[draw,text width=36mm,align=center] at (0,-1.3) {基线已有结局$\to$排除};
\node[draw=dcrgold,text width=25mm,align=center] (cases) at (8,1.2) {\textcolor{dcrgold}{病例}\\测量预测因素};
\node[draw=dcrgold,text width=25mm,align=center] (controls) at (8,-.2) {\textcolor{dcrgold}{对照}\\测量预测因素};
\node[draw=dcrgold,text width=25mm,align=center] at (8,-1.8) {队列其余成员\\多数无须测量\\预测因素};
\draw[->,thick,dcrteal] (2,.7)--(6.6,.7);
\node[text=dcrgold,align=center] at (4.4,1.4) {结局发生时或\\随访结束时测量};
\foreach \x in {3.5,5.5}{\draw[->,dcrteal] (\x,.7) to[bend left=30] (\x+.5,-2.1);}
\node[text=dcrteal] at (4.7,-2.5) {失访};
\draw[->,thick,dcrteal] (2,-3)--(6.5,-3);\node[fill=white,text=dcrteal] at (4.4,-3) {随访时间};
\node[align=center] at (0,-3) {队列在过去组建};
\node[align=center] at (8,-3.65) {现在从保存标本、影像等\\测量预测变量};
\node[align=left] at (10.4,1.2) {全部病例\\或其样本};
\node[align=left] at (10.4,-.2) {对照样本};
\end{tikzpicture}}
% Each of the 16 timelines below follows the printed Figure 9.4 (PDF p258).
% The three incident cases, two losses, and three changes in exposure are retained.
\providecommand{\chNineRiskSets}{\begin{tikzpicture}[font=\footnotesize,>=Stealth,x=.019cm,y=-.019cm]
\draw[dashed,dcrred] (171,207) ellipse[x radius=3.12cm,y radius=3.17cm];
\node[text=dcrred,align=left,anchor=west] (pop) at (8,16) {可接触人群};
\draw[->,dcrred] (pop.east) to[bend left=20] (126,48);
\node[align=center,anchor=south] at (177,104) {符合选择标准的队列};
\draw[->] (225,94) to[bend left=20] (242,111);
\draw[dcrred] (87,111) rectangle (273,299);
\node[align=center,text width=29mm,font=\footnotesize\bfseries] at (180,205) {有结局发生风险的\\合格参与者};
% 1: exposure starts before the first case.
\draw[dcrblue,thick] (273,111)--(325,111);
\draw[red,thick] (325,111)--(355,111);
\node[diamond,fill=red,inner sep=1.7pt] at (355,111) {};
% 2--7: one loss and a change from exposed to unexposed person-time.
\draw[->,dcrblue,thick] (273,124)--(377,124);
\draw[->,dcrblue,thick] (273,136)--(389,136);
\draw[dcrblue,thick] (273,149)--(389,149);
\draw[dcrblue,thick] (389,144)--(389,154);
\draw[->,dcrblue,thick] (273,161)--(407,161);
\draw[red,thick] (273,173)--(378,173);
\draw[->,dcrblue,thick] (378,173)--(418,173);
\draw[->,dcrblue,thick] (273,186)--(438,186);
% 8: the second case occurs during unexposed person-time.
\draw[dcrblue,thick] (273,199)--(439,199);
\node[diamond,fill=dcrblue,inner sep=1.7pt] at (439,199) {};
% 9--12: the second loss is on the eleventh timeline.
\draw[->,red,thick] (273,211)--(466,211);
\draw[->,dcrblue,thick] (273,224)--(479,224);
\draw[dcrblue,thick] (273,236)--(494,236);
\draw[dcrblue,thick] (494,231)--(494,241);
\draw[->,dcrblue,thick] (273,248)--(504,248);
% 13: the third case occurs during exposed person-time.
\draw[red,thick] (273,260)--(507,260);
\node[diamond,fill=red,inner sep=1.7pt] at (507,260) {};
% 14--16: the final participant becomes exposed during follow-up.
\draw[->,dcrblue,thick] (273,273)--(528,273);
\draw[->,dcrblue,thick] (273,285)--(541,285);
\draw[dcrblue,thick] (273,298)--(406,298);
\draw[->,red,thick] (406,298)--(553,298);
\foreach \xx/\nn/\labelx in {357/1/310,439/2/420,508/3/530}{
  \draw[dcrteal] (\xx-10,98) rectangle (\xx+10,313);
  \node[draw=dcrteal] (r\nn) at (\labelx,374) {风险集\nn};
  \draw[->,dcrteal] (\xx,313)--(r\nn.north);
}
% Legend preserves the meanings of colors and endpoint symbols.
\draw[dcrblue,thick] (540,89)--(568,89);
\node[anchor=west] at (575,89) {未暴露人时};
\draw[red,thick] (540,116)--(568,116);
\node[anchor=west] at (575,116) {暴露人时};
\node[diamond,fill=red,inner sep=1.6pt] at (545,143) {};
\node[diamond,fill=dcrblue,inner sep=1.6pt] at (562,143) {};
\node[anchor=west] at (575,143) {新发病例};
\draw[red,thick] (547,165)--(547,178);
\draw[dcrblue,thick] (562,165)--(562,178);
\node[anchor=west] at (575,172) {失访};
\draw[->,dcrblue,thick] (540,205)--(568,205);
\node[anchor=west,align=left] at (575,213) {随访结束时\\未成为病例};
\end{tikzpicture}}

\begin{figure}[htbp]\centering\chNineCaseControl
\caption*{\textbf{图9.1\quad 病例对照研究。}界定选择标准，从病例人群抽一个样本，从对照人群抽另一个样本；测量相关变量当前值，通常辅以既往资料。}
\end{figure}

\begin{bookbox}{例9.1\quad 疾病暴发的病例对照研究}
2018年澳大利亚多州发生甲型肝炎（Hepatitis A）暴发，引发公共卫生部门调查（2）。病例为2018年4月13日至6月8日确认的IB基因型甲肝患者；对照患其他法定报告传染病，如沙门菌病（salmonellosis）、弯曲菌病（campylobacteriosis）或隐孢子虫病（cryptosporidiosis）。\sourcepage{245}对照来自病例所在或相邻地方政府辖区，按年龄组（0～14、15～39、至少40岁）作频数匹配（frequency matching），比例2:1，即每一年龄组每例病例选两名同组对照（见第10章）。两组均被询问发病前可能的暴露。

与病例身份关联最强的是食用冷冻石榴籽：13例病例中9人，而21名对照中仅1人，比值比（OR）45.0，95\%置信区间3.8～2065（见附录9A）。此后，从一家埃及厂商的两包冷冻石榴籽检出甲肝病毒；该厂商产品曾与2012年不列颠哥伦比亚省甲肝暴发相关（3）。澳大利亚农业与水资源部因此开始检查、检测该厂商此后所有到货批次。
\end{bookbox}

若把研究设计比作酒单，病例对照就是“店家红酒”：比其他选择朴素、风险略高，却便宜得多，有时出奇地好。因为偏倚机会更多，设计颇具挑战，但良好设计产生重要结果的例子很多，如母亲使用己烯雌酚（diethylstilbestrol）与女儿阴道癌（仅七例就得出明确结论的经典研究！）（4），以及俯卧睡眠与婴儿猝死综合征（5）；后者的简单发现挽救了数千生命（6）。

病例对照不能估计疾病发病或患病率，因为研究中患病者比例由抽取多少病例、多少对照决定，而非人群频率。不过，它能描述病例特征，更重要的是，估计每个预测变量与结局关联的强度。\sourcepage{246}该估计以比值比表示；若暴露与未暴露者结局风险均低（约10\%以下），它近似风险比（见附录9B）。

病例对照最初用于寻找疾病危险因素，所以为方便叙述，我们通常称患病者为“病例”，并称两组间比较的暴露为“危险因素”。但它也能研究其他少见结局，如通常不致死疾病患者的死亡。有时关注的暴露被假设可预防疾病，如疫苗或治疗效果的病例对照研究；这时研究者希望显示暴露组风险较低，即比值比小于1。

\subsection{病例对照研究的优点}
\subsubsection{研究罕见结局效率高}
一大优点是用较少参与者快速取得信息。以包皮环切对随后阴茎癌的作用为例：此癌在环切男性中极少，在未环切者也少见，后者终生累积发病率约0.16\%（7）。若要以合理把握（80\%）检出环切是否使阴茎癌风险增加50倍，假定约一半接受环切，队列需随访6000多名男性数十年。\footnote{译注：底本此处原文为“circumcision increases the risk … 50-fold”（环切使风险增加50倍），与前文称环切者更少患癌的叙述不一致；此处保留原文方向。}出生时环切的随机试验需同样样本，但病例发生的中位时间在入组67年后，需要两三代研究者接续随访！

对同一问题，病例对照以同样把握检出同样相对风险，只需16例病例与16名对照。当然仍须找到有代表性的16例；病越罕见，寻找可能越费力，见下文缺点。\sourcepage{247}但对罕见病，或暴露与疾病间潜伏期很长者，病例对照不仅效率远高于其他设计，常是唯一可行选择。

\subsubsection{有助于提出假设}
回顾性方法和可考察多个预测变量，使病例对照适合提出新暴发病因的假设。暴发通常由传染病引起（例9.1），但并非总如此。例如，海地儿童急性肾衰竭死亡流行的病例对照研究（8），发现服用当地生产的对乙酰氨基酚糖浆的比值比为53。后续确认肾衰竭源于污染物二甘醇（diethylene glycol）中毒；遗憾的是，这类问题至今仍发生（9、10）。

\subsection{病例对照研究的缺点}
它也有重要缺点。第一，只能研究一种结局，即作为两组抽样标准的疾病有无；队列、横断面和临床试验可研究多个结局。第二，信息有限：除非另知病例来源人群和时间段，不能直接估计疾病发病或患病率。最大的弱点是易受偏倚影响，主要来自病例和对照分别选择，以及预测变量回顾性测量。下面分别讨论及其应对方法。

\subsection{选择偏倚及其避免方法}
\subsubsection{寻找与抽取病例}
抽样从病例开始。理想样本包括明确人群内全部新发病例，或其随机样本。\sourcepage{248}但立刻遇到问题：如何知道谁患病、谁没有？横断面和队列对全部参与者系统确定疾病；病例对照则只能从已经确诊且可接受研究的患者中抽样。未诊断、误诊、无法参加或已死亡者难以纳入，因此样本可能不能代表全部患病者（图9.2）。

\begin{figure}[htbp]\centering\chNineSelection
\caption*{\textbf{图9.2\quad 病例对照研究中的病例可能不能代表全部新发病例的若干原因。}}
\end{figure}

一般来说，当病例样本在\emph{所研究危险因素方面}不具代表性时，选择偏倚才重要。髋部骨折、创伤性截肢等几乎总须住院且诊断确认简单的疾病，可从确诊且可接受研究者中抽样，至少在就医便利的人群如此。巢式病例对照的病例来自明确队列；若随访足够长，疾病最终几乎总能诊断，则病例也可近似全体新发病者。

\sourcepage{249}但并非总会就诊的病不易用病例对照研究，因为诊断前已有选择。例如，有研究担心输精管结扎增加前列腺癌风险（11），但诊断可能依赖筛查，因此病例可能过多代表接受过医疗、特别是泌尿科照护的男性，如接受结扎者。事实上，结扎者更常接受前列腺癌筛查（12）。

\subsubsection{抽取对照}
必须思考上述问题，但病例来源常受可及性限制，样本未必完全有代表性，却可能是研究者所能得到的全部。更难、选择空间更大的决定，是如何选合适对照。总目标是从这样的人群抽取：\emph{如果他们患病，就会成为本研究病例}。有以下三种策略。

\paragraph{门诊、医院或登记来源对照。}最方便的对照常来自病例就诊的同一门诊、医院或同一登记，因为容易接触。医院／门诊对照还可能部分补偿同来源病例的选择偏倚，因为他们若患研究疾病，很可能也会到此就诊。例9.1的甲肝调查中，其他法定报告传染病患者的联系资料易得，而且假如患甲肝，推测也会被报告为病例。

\sourcepage{250}但用不具代表性的对照补偿不具代表性的病例，可能有问题。若所研究危险因素引起对照求医的疾病，对照中该因素频率就会虚高，削弱甚至反转其与结局的关联。若甲肝研究关注的是一般食品和饮水安全相关暴露，以沙门菌、弯曲菌、隐孢子虫患者为对照就不妥，因为食物卫生差者过多。又如，一项将儿童、青年居所靠近油气井与血液系统肿瘤联系的研究（13），使用登记中非血液系统肿瘤者为对照；若油气井周边化学物也引起这些其他癌症，血液肿瘤关联就会被低估。

其他疾病对照所患疾病可能也是所研究危险因素所致，造成误导。因此必须权衡方便性是否值得承担对效度的威胁。

\paragraph{使用基于人群的病例与对照样本。}疾病登记，以及地区人群和健康计划的临床、行政数字资料迅速增加，使许多疾病可开展基于人群的病例对照研究。此类病例通常代表地区或健康计划中患病人群，从而简化对照选择：可从登记覆盖人群或数据库中抽取有代表性的“非病例”。

有登记等数据库时，基于人群的病例对照是最理想的设计。登记越完整，覆盖人群越固定、迁入迁出越少，就越接近巢于队列或临床试验的病例对照，\sourcepage{251}前提是能找到并招入对照。若人群已列册，且病历可供研究者使用，后两项任务较容易。

但即使卫生系统整合良好，图9.2所示识别病例的筛选过程，对不同人群仍可能影响不同。只要诊断须进入医疗系统，或纳入研究须联系及同意，就可能引入偏倚，因为不说英语、较不信任科学机构等人可能更少入选。

\paragraph{使用两个以上对照组。}选择对照很棘手，尤其病例不具代表性时，有时应按不同方式选两个以上对照组。例如，海地口服霍乱疫苗效果研究（14），比较病例与两组对照接种史。“检测阴性”对照在同一霍乱治疗中心就诊，有水样腹泻，但霍乱检测阴性；另一“社区”组无腹泻，每例最多四名，按年龄组、时间和居住邻里匹配。自报两剂疫苗的效果，两个对照组估计相近，分别73\%和74\%，令人放心。

遗憾的是，不同对照选择策略的偏倚有时使结果冲突，暴露出该问题上病例对照设计的脆弱性。此时应搜集更多信息，确定每组潜在偏倚大小（第10章）。\sourcepage{252}例如，担心门诊对照的吸烟过多，是因为吸烟促使他们因疲乏等其他问题就诊，可考察对照中危险因素频率随就诊原因怎样变化。无论如何，结果不一致而承认答案未知，总比只有一组对照却得出错误结论好。

上述策略均可配合\term{匹配}（matching），使病例和对照在与疾病有关或可能有关、但非主要关注暴露的因素上可比。例如，许多危险因素、疾病与年龄、性别、地区相关；若两组这些变量可比，结果更可信。可选择与病例年龄、性别及某种地区指标匹配的对照。但匹配也有缺点，尤其匹配收入、血清胆固醇等可改变的预测变量时。原因与通常更好的替代方法见第10章。

\subsection{差异性测量偏倚及其避免方法}
第二大弱点是测量误差所致偏倚风险增加，来自回顾性测量暴露；测量者可能是研究者、临床医师（病历记录），或回忆多年前暴露的病例、对照本人。人们记忆不完整；若病例和对照同样不完整，称为暴露的\term{非差异性错分}（nondifferential misclassification），会使关联更难发现，流行病学上即比值比趋向1。更令人担心的是，患病确诊可能使病例与对照以不同方式回忆或报告暴露；\sourcepage{253}这种差异性暴露错分称为\term{回忆偏倚}（recall bias），对测得关联的影响不可预测。

例如，日晒与恶性黑色素瘤关系被广泛宣传，可能使确诊患者不同于对照地回忆日晒史。Cockburn等（15）巧妙研究一人患黑色素瘤、一人未患的双生子；问患病者童年谁更多日光浴，匹配比值比为2.2（95\%置信区间1.0～4.7），问未患病同胞则仅0.8（0.4～1.8）。但谁更容易晒黑、晒伤等其他问题，未见回忆偏倚。

队列在诊断前询问暴露，因此不会有这种回忆偏倚。另一项巢于数年前已收集日晒资料队列中的黑色素瘤病例对照，直接检验了回忆偏倚：比较诊断前后两组自报日晒（16），发现两组回忆都有不准，但回忆偏倚很少。因此应考虑可能性，却并非必然发生（17）。

除第4章的测量偏倚控制方法——标准化变量操作定义、选择客观方法、关键变量用多来源资料补充等——还有两种专门策略：
\begin{itemize}
\item \textbf{使用结局发生前记录的资料。}如前述霍乱疫苗效果研究（14），除自报外也用实际接种记录。但该好方法受危险因素记录可得性和可信度限制；霍乱研究中，自报接种的病例和对照都只有约一半有验证记录。\sourcepage{254}
\item \textbf{使用盲法。}一般方法见第4章，病例对照访谈还有特定问题。理论上，观察者和参与者都可不知病例／对照身份及所研究危险因素，因此有四种可能（表9.1）。
\end{itemize}

\begin{table}[htbp]\centering\small
\caption*{\textbf{表9.1\quad 病例对照研究的盲法}}
\begin{tabularx}{\textwidth}{L{22mm}Y Y}\toprule
设盲对象&对病例／对照身份设盲&对危险因素测量设盲\\\midrule
参与者&若病例和对照所患疾病均可能与危险因素相关，则可行&加入“假”危险因素；若两组有差异应警觉。若危险因素已广为宣传，可能无效\\
观察者&若病例外表无法与对照区分则可行，但细微信号及参与者主动透露的信息会使其困难&访谈者并非研究者时可行，但可能难以维持\\\bottomrule
\end{tabularx}
\end{table}

理想状态是参与者和观察者都不知谁是病例或对照，实际常很难。参与者知道自己是否生病，只有对照也患某种自认为与所研究危险因素可能相关的疾病，才可能对身份设盲。某些疾病外观明显（黄疸、喉切除难不被注意），访谈者还可从回答发现线索。

对具体危险因素设盲通常较容易。病例对照常是调查疾病的第一步，未必只关注一个因素。\sourcepage{255}可加入看似合理但不相关的“假”问题，使双方不了解研究假设。例如，研究蜂蜜与婴儿肉毒中毒，可同样详细询问酸奶、香蕉。这不能阻止差异性偏倚，却可估计是否存在：若病例只报告更多蜂蜜，而其他食物不增，差异性测量偏倚较不可能。但若蜂蜜与肉毒中毒关联已被广泛宣传，或假因素竟是真因素，此策略便无效。

对血液检查、X线等测量，观察者不知道病例／对照身份尤为有效，且容易、应始终做到：由非测量者给标本（或患者）贴编码标签即可。15项比较髋部骨折患者与对照骨量的研究显示，未设盲者发现的差异远大于设盲者（18）。

\section{巢式病例对照、发病密度巢式病例对照与病例队列研究}
巢式病例对照就是在明确队列内“嵌入”病例对照研究（图9.3）。队列可能早已作为正式队列研究的一部分建立，并存储标本、影像等，等待结局发生后分析；也可在尚未界定的队列中从头设计，此时第一步是描述队列。
\sourcepage{256}
\begin{figure}[htbp]\centering\chNineNested
\caption*{\textbf{图9.3\quad 巢式病例对照研究。}确定已保存标本、影像等资料的队列；测量区分病例与对照的结局变量；对全部病例和部分对照，在队列建立时保存的标本、影像等中测量预测变量，并测量其他变量。}
\end{figure}

首先找出对结局处于风险中的队列，足够大以产生回答问题所需的病例数，并能测量暴露——有保存标本，或有暴露信息的病历。按第8章方法明确纳入、排除标准，界定风险人群。还必须明确每人入队列日期，可是固定日（如2021年1月1日符合标准的所有健康计划成员），也可是每人风险期开始的不同日期（如入队列日，\sourcepage{257}或再发心肌梗死危险因素研究中的首次梗死日）。

接着描述结局发生标准，要求在入队列后、规定随访结束前出现。若结局罕见、随访近乎完整、一次基线暴露测量足够，则很简单：找出随访结束前全部发生结局者作为病例，从同队列未发生者中随机抽对照，再测量两组预测因素，比较危险因素水平。

若随访时间不同或不完整，或暴露随时间改变，只在病例与随机对照入队列时测一次暴露就不足。此时宜设计发病密度巢式病例对照，从\term{风险集}（risk sets）抽对照。每当一例病例发生，风险集包括已随访与该病例同样长时间、但尚未成为病例的队列成员（图9.4）。与其他病例对照匹配一样，按随访时间匹配也必须在分析中考虑。
\sourcepage{258}
\begin{figure}[htbp]\centering\chNineRiskSets
\caption*{\textbf{图9.4\quad 发病密度巢式病例对照研究。}可为前瞻性或回顾性。前瞻性步骤：界定选择标准并招入队列；确定每人的入队列日期，使随访时间有共同起点；保存标本、影像等；随访并确定病例及确诊日；每例病例从已随访相同时长、在确诊时尚未发病、死亡或失访者组成的风险集中抽一个以上对照；在病例及匹配对照自基线保存的标本、影像等中测量预测变量及其他当前变量。回顾性设计的前四步已经完成。}
\end{figure}

若入队列是固定日（如2018年1月1日），2019年7月1日确诊病例的对照，就从该日尚未发生结局者中抽取。若入队列日期不同，入组18个月后确诊者的对照，就从随访18个月尚未发病者中抽。\sourcepage{259}根据对致病暴露时间及持续长度的假设，可比较入组暴露值、平均值，或确诊前固定时点（如3个月）的值。

按风险集抽样意味着，某人可先成为早发病例的对照，后来自己成为病例，其间暴露值甚至会改变。实质上，在适当统计分析帮助下，设计以病例发生时点划分一段段风险人时，用该时段预测变量值（有时也用前一时段，即\term{滞后变量}，lagged variables）预测该人时段的病例发生。这称为发病密度设计（例9.2）。

\begin{bookbox}{例9.2\quad 意大利阿片过量死亡的发病密度巢式病例对照研究}
阿片过量是高收入国家主要死因之一，美沙酮等阿片激动剂治疗可能降低风险。使用物质者会进出治疗，治疗随时间变化；意大利、澳大利亚和英国研究者因而采用发病密度巢式设计，高效量化阿片激动剂及其他治疗的益处（19）。
\begin{enumerate}
\item \textbf{确定队列与风险期。}研究巢于VEdeTTE队列，1998—1999年从意大利公共治疗中心招募海洛因使用者。确定两个地区4444名参与者自1998年9月至2005年12月31日的死亡。
\item \textbf{确定病例及发生日期。}不知道参与者病例／对照身份的工作人员从临床记录提取生死状态。\sourcepage{260}随访期死亡316人，其中95人编码为过量死亡。
\item \textbf{从每例匹配风险集抽对照。}每个致死过量病例，随机选择四名其死亡当日（对照的索引日）仍存活者，匹配地区、年龄（$\pm5$岁）和性别。一名对照可匹配多个病例，后来的病例也可给更早死亡病例作对照。
\item \textbf{测量病例与对照的预测因素。}研究者在不知身份时，从临床记录提取病例死前两个月、对照索引日前两个月的药物治疗类型。两个月均未治疗者，提取末次治疗类型与日期。相比队列，发病密度设计只需在少得多的人、短得多的时段提取这些资料。
\end{enumerate}
研究者恰当地使用\term{条件logistic回归}（conditional logistic regression）分析资料，发现阿片激动剂治疗与死亡比值降低90\%以上有关（OR＝0.09，95\%置信区间0.03～0.24）；调整无家可归、HIV阳性、饮酒、法律问题和基线报告的过量后未改变结果（调整OR＝0.08，0.03～0.23）。
\end{bookbox}

巢式病例队列与巢式病例对照相似，但不是从未发生结局者抽对照，而是从全部队列成员随机抽样，不考虑结局。随机子队列内少数人可能已发生所研究结局；结局罕见时这个数很小。优势是同一个随机样本可为多个不同结局提供对照，\sourcepage{261}如例9.3，且可提供整个队列危险因素总体流行程度。

\begin{bookbox}{例9.3\quad 血浆ACE2与死亡或心脏代谢疾病风险：巢式病例队列研究}
循环中的血管紧张素转换酶2（angiotensin-converting enzyme 2，ACE2）水平，可能是肾素－血管紧张素系统失调的标志。主要来自加拿大安大略省汉密尔顿人口健康研究所的研究者，考察血浆ACE2水平对数种不良健康结局的预测能力（20）。
\begin{enumerate}
\item \textbf{确定队列。}采用前瞻性城乡流行病学研究（Prospective Urban Rural Epidemiology，PURE），其基线纳入27个低、中、高收入国家35～70岁人群，并将14国部分参与者基线血样存于$-165\,^{\circ}\mathrm{C}$。有可分析保存样本、且自报为居住国主要族群者合格（如瑞典的欧洲血统）。样本2005—2006年采集，中位随访9.4年。
\item \textbf{确定病例。}结局事件为死亡（$N=1985$）、心肌梗死（882）、卒中（663）、心力衰竭（264）和糖尿病（1715），已由PURE确定。
\item \textbf{从队列抽样。}从有可分析样本的55246人随机抽5084人组成“子队列”，不因发生所研究事件排除，因此是病例队列而非巢式病例对照。
\item \textbf{测量病例与子队列预测因素。}在生物样本库标本测ACE2，分析先前收集的性别、体质指数、吸烟、血压等传统心血管危险因素。\sourcepage{262}
\end{enumerate}
ACE2对全部结局均有很强预测能力；调整传统心血管因素后，每升高一个标准差（21），风险率比为1.21～1.44，其预测强于吸烟、糖尿病、血压、血脂和体质指数等临床危险因素。
\end{bookbox}

\subsection{优点}
巢式病例对照与病例队列尤其适合研究开始时存档、留待日后分析的血清、其他标本或影像的昂贵测量。只对全部病例和部分对照测量，比测整个队列便宜得多。测量者设盲时，保留了队列在结局前收集预测变量的优点；并避免传统病例对照无法测量死亡病例，以及从不同人群抽病例、对照的潜在偏倚。

\subsection{缺点}
仍与其他观察性设计一样，关联可能由未测量或测量不精密的混杂因素造成，基线测量也可能受无症状临床前疾病影响，即“果因”而非因果（第10章）。

\subsection{其他考虑}
这些设计本应更广使用。计划大型前瞻性研究时，应考虑保存生物样本（如冻存血清库），或储存分析昂贵的影像、记录，供以后巢式分析。\sourcepage{263}须确保储存条件能长期保存关注物质。随访中收集新样本或新信息也有用，可在后来的发病密度病例对照中高效使用。

\section{病例交叉研究}
病例交叉是病例对照的变体，适用于间歇性暴露的短期效应。回顾性研究从已经有结局的病例开始；与传统病例和另一组对照比较不同，每个病例为自己的对照。比较其结局发生时（或紧前）的暴露，与自身一个以上其他时点暴露。

例如，McEvoy等研究有手机或使用手机的车祸伤者（22），用电话公司记录，比较事故前10分钟与24小时、72小时、7天前同一时刻驾驶时的手机使用，发现车祸前10分钟使用可能性约为对照时段的四倍。

分析类似匹配病例对照，只是对照暴露来自同一病例不同时段，而非匹配对照；见附录9A第三种情形。大型人群中也用病例交叉研究颗粒物空气污染等时间变化暴露，发现其与住院（23）、老年人全因死亡（24）、院外心脏骤停（25），甚至婴儿死亡（26）相关。

\section{选择观察性设计}
\sourcepage{264}
表9.2总结两章主要设计的优缺点。详细讨论已在前文，此处只强调：没有一种最好或最差；每种都有其位置和用途，取决于研究问题与具体条件。

\begin{longtable}{L{30mm}L{58mm}L{61mm}}
\caption*{\textbf{表9.2\quad 主要观察性设计的优缺点}}\\\toprule
设计&优点&缺点$^{a}$\\\midrule\endfirsthead
\toprule 设计&优点&缺点$^{a}$\\\midrule\endhead
横断面&持续时间较短；队列或临床试验的良好起点；提供多个预测因素和结局的患病率&不能确定事件顺序；罕见预测因素或罕见结局不可行；不能给出发病\\
队列：共同特征&确定事件顺序；多个预测因素和结局；事件数随时间增加；给出发病、相对风险、超额风险&常需大样本；罕见结局较不可行\\
前瞻性队列&更能控制参与者选择与测量；避免预测因素测量偏倚&随访可能很长；常昂贵\\
回顾性队列&随访已经完成；相对便宜&对选择和测量控制少\\
多队列&不同队列有不同或罕见暴露时有用&从不同人群抽样造成偏倚与混杂\\
\sourcepage{265}病例对照&适合罕见结局；时间短、样本小；相对便宜&两个来源人群抽样的偏倚与混杂；差异性测量偏倚；限一种结局；事件顺序可能不清；除非巢于队列，否则无患病率、发病或超额风险\\
巢式病例对照&具有回顾性队列优点；预测因素测量昂贵时费用更低&若危险因素未曾测量，也非来自已存标本或影像，测量仍易偏倚；通常需既有明确队列\\
发病密度巢式病例对照&分析风险关系时可考虑危险因素水平随时间的变化与失访&需随访中各时点危险因素水平和病例发生的测量；通常需既有明确队列\\
巢式病例队列&同巢式病例对照；且一个对照组可用于多个不同结局研究&同巢式病例对照\\
病例交叉&病例自为对照，减少随机误差与混杂&要求暴露只具有即时、短期效应\\\bottomrule
\multicolumn{3}{p{154mm}}{\footnotesize $^{a}$相对于随机试验，全部观察性设计均易受混杂变量影响，见第10章。}\\
\end{longtable}

\section{小结}
\sourcepage{266}
\begin{enumerate}
\item 病例对照比较有结局病例样本与无结局对照样本中危险因素的频率。分别抽取有病、无病者，费用较低，研究罕见病效率特别高。
\item 易受选择偏倚影响。可让病例、对照按同样（虽不具代表性）方式抽样，开展基于人群研究，或使用多种方式抽取多个对照组；均可配合年龄、性别、地区等匹配。
\item 回顾性设计还易产生差异性测量偏倚，可用结局前预测因素测量，并对参与者与工作人员设盲。
\item 避免两类偏倚的最好方法是巢式病例对照，从较大队列中随机抽病例和对照；血清、影像等昂贵测量也只需在研究结束后对少数人做。
\item 发病密度设计高效分析风险关系，可考虑危险因素水平及可随访性随时间变化。
\item 巢式病例队列用全队列随机样本代替非病例，可作多个结局的对照，并直接提供危险因素总体流行程度。
\item 病例交叉是匹配病例对照的变体，用两个以上时点的观察，让病例作自身对照。
\end{enumerate}
\begin{references}
\sourcepage{267}
\item Guenel P, Laforest L, Cyr D, et al. Occupational risk factors, ultraviolet radiation, and ocular melanoma: a case-control study in France. Cancer Causes Control. 2001;12(5):451-459.
\item Franklin N, Camphor H, Wright R, Stafford R, Glasgow K, Sheppeard V. Outbreak of hepatitis A genotype IB in Australia associated with imported frozen pomegranate arils. Epidemiol Infect. 2019;147:e74.
\item Swinkels HM, Kuo M, Embree G, et al. Hepatitis A outbreak in British Columbia, Canada: the roles of established surveillance, consumer loyalty cards and collaboration, February to May 2012. Euro Surveill. 2014;19(18):20792.
\item Herbst AL, Ulfelder H, Poskanzer DC. Adenocarcinoma of the vagina. Association of maternal stilbestrol therapy with tumor appearance in young women. N Engl J Med. 1971;284(15):878-881.
\item Beal SM, Finch CF. An overview of retrospective case-control studies investigating the relationship between prone sleeping position and SIDS. J Paediatr Child Health. 1991;27(6):334-339.
\item Mitchell EA, Blair PS. SIDS prevention: 3000 lives saved but we can do better. N Z Med J. 2012;125(1359):50-57.
\item Kochen M, McCurdy S. Circumcision and the risk of cancer of the penis. A life-table analysis. Am J Dis Child. 1980;134(5):484-486.
\item O'Brien KL, Selanikio JD, Hecdivert C, et al. Epidemic of pediatric deaths from acute renal failure caused by diethylene glycol poisoning. Acute Renal Failure Investigation Team. JAMA. 1998;279(15):1175-1180.
\item Fatal poisoning among young children from diethylene glycol-contaminated acetaminophen—Nigeria, 2008–2009. MMWR Morb Mortal Wkly Rep. 2009;58(48):1345-1347.
\item Llamas M. Drug Makers Warned for Potential Diethylene Glycol Toxin Contamination: Drugwatch.com; 2020 [updated April 20, 2020]. \url{https://www.drugwatch.com/news/2020/04/20/diethylene-glycol-toxin-contamination/}
\item Nutt M, Reed Z, Kohler TS. Vasectomy and prostate cancer risk: a historical synopsis of undulating false causality. Res Rep Urol. 2016;8:85-93.
\item Shang Y, Han G, Li J, et al. Vasectomy and prostate cancer risk: a meta-analysis of cohort studies. Sci Rep. 2015;5:9920.
\item McKenzie LM, Allshouse WB, Byers TE, Bedrick EJ, Serdar B, Adgate JL. Childhood hematologic cancer and residential proximity to oil and gas development. PLoS One. 2017;12(2):e0170423.
\item Franke MF, Jerome JG, Matias WR, et al. Comparison of two control groups for estimation of oral cholera vaccine effectiveness using a case-control study design. Vaccine. 2017;35(43):5819-5827.
\sourcepage{268}
\item Cockburn M, Hamilton A, Mack T. Recall bias in self-reported melanoma risk factors. Am J Epidemiol. 2001;153(10):1021-1026.
\item Parr CL, Hjartaker A, Laake P, Lund E, Veierod MB. Recall bias in melanoma risk factors and measurement error effects: a nested case-control study within the Norwegian Women and Cancer Study. Am J Epidemiol. 2009;169(3):257-266.
\item Gefeller O. Invited commentary: recall bias in melanoma—much ado about almost nothing? Am J Epidemiol. 2009;169(3):267-270; discussion 71-72.
\item Cummings SR. Are patients with hip fractures more osteoporotic? Review of the evidence. Am J Med. 1985;78(3):487-494.
\item Faggiano F, Mathis F, Diecidue R, et al. Opioid overdose risk during and after drug treatment for heroin dependence: an incidence density case-control study nested in the VEdeTTE cohort. Drug Alcohol Rev. 2021;40(2):281-286.
\item Narula S, Yusuf S, Chong M, et al. Plasma ACE2 and risk of death or cardiometabolic diseases: a case-cohort analysis. Lancet. 2020;396(10256):968-976.
\item Newman TB, Browner WS. In defense of standardized regression coefficients. Epidemiology. 1991;2(5):383-386.
\item McEvoy SP, Stevenson MR, McCartt AT, et al. Role of mobile phones in motor vehicle crashes resulting in hospital attendance: a case-crossover study. BMJ. 2005;331(7514):428.
\item Wei Y, Wang Y, Di Q, et al. Short term exposure to fine particulate matter and hospital admission risks and costs in the Medicare population: time stratified, case crossover study. BMJ. 2019;367:l6258.
\item Di Q, Dai L, Wang Y, et al. Association of short-term exposure to air pollution with mortality in older adults. JAMA. 2017;318(24):2446-2456.
\item Kojima S, Michikawa T, Matsui K, et al. Association of fine particulate matter exposure with bystander-witnessed out-of-hospital cardiac arrest of cardiac origin in Japan. JAMA Netw Open. 2020;3(4):e203043.
\item Scheers H, Mwalili SM, Faes C, Fierens F, Nemery B, Nawrot TS. Does air pollution trigger infant mortality in Western Europe? A case-crossover study. Environ Health Perspect. 2011;119(7):1017-1022.
\end{references}


\chapterappendix{9A}{病例对照研究比值比的计算}
\sourcepage{269}
\subsection{病例对照研究}
例9.1调查甲肝暴发：13例病例中9人、21名对照仅1人报告吃过冷冻石榴籽（表9A.1）。比值比为病例暴露比值$a/c$除以对照暴露比值$b/d$，数学上等于$ad/bc$。
\begin{table}[htbp]\centering\small
\caption*{\textbf{表9A.1\quad 例9.1的比值比计算}}
\begin{tabular}{lrr}\toprule
&\multicolumn{2}{c}{结局变量：诊断甲肝}\\
预测变量：吃过冷冻石榴籽&是&否\\\midrule
是&9（$a$）&1（$b$）\\
否&4（$c$）&20（$d$）\\
合计&13&21\\\bottomrule
\end{tabular}
\end{table}
疾病罕见（约低于10\%）时，如本例甲肝，比值比很好地估计风险比。因此，食用污染石榴籽者患甲肝可能性约为未食用者45倍。
\[
\text{相对风险}\approx\text{比值比}=\frac{a/c}{b/d}=\frac{ad}{bc}=\frac{9\times20}{4\times1}=45.
\]

\subsection{匹配病例对照研究}
为说明与病例交叉分析的相似性，使用同一问题：手机使用是否增加手机拥有者车祸风险？传统匹配病例对照可把自报驾驶时用手机的频率作为危险因素，以事故者为病例，未事故者为对照，按年龄、性别及手机号码前缀匹配，\sourcepage{270}询问是否曾在驾驶时用手机。为简单，将暴露二分为使用者与非使用者。每对依次分为：两人均使用、均不使用、仅病例使用、仅对照使用。假如有300对，结果可如表9A.2。

\begin{table}[htbp]\centering\small
\caption*{\textbf{表9A.2\quad 假想配对病例对照研究：驾驶时习惯性使用手机与车祸伤}}
\begin{tabular}{lrrr}\toprule
&\multicolumn{3}{c}{病例（有车祸伤）}\\
匹配对照&使用者&非使用者&合计\\\midrule
使用者&110&40&150\\
非使用者&90&60&150\\
合计&200&100&300\\\bottomrule
\end{tabular}
\end{table}
90对是病例使用、对照不用，40对相反。与甲肝表每格为\emph{人数}不同，匹配研究的每格为参与者\emph{对数}；故表9A.2共有600人，即300病例、300对照。

其比值比只取“病例暴露、对照不暴露”对数，除以“对照暴露、病例不暴露”对数，即\term{不一致对}（discordant pairs）之比。\sourcepage{271}暴露相同的\term{一致对}（concordant pairs）不贡献暴露－结局关联信息。表中OR＝90／40＝2.25；车祸罕见，故近似风险比。假如结果真实，使用手机者的车祸风险超过不用者两倍。

\subsection{病例交叉研究}
同问题的病例交叉资料来自McEvoy等（22）（表9A.3）。每格为病例人数，而非配对数，但代表同一病例两个时段：事故紧前与7天前比较时段。左上5人表示两时段都用手机；其下27人表示事故前用，7天前不用。

\begin{table}[htbp]\centering\small
\caption*{\textbf{表9A.3\quad 病例交叉研究：近期手机使用与车祸伤的比值比}}
\begin{tabular}{lrrr}\toprule
&\multicolumn{3}{c}{事故时段}\\
事故前7天&司机使用手机&未使用&合计\\\midrule
司机使用手机&5&6&11\\
未使用&27&288&315\\
合计&32&294&326\\\bottomrule
\end{tabular}
\end{table}
另有6人在事故时不用、7天前用。\sourcepage{272}比值比为不一致时段数之比，27／6＝4.5，即驾车用手机时段的车祸风险为不用时段的4.5倍。

\chapterappendix{9B}{比值比何时、为何近似风险比}
\sourcepage{273}
暴露与疾病的关联可用四格表表示：
\begin{center}\begin{tabular}{lcc}\toprule&患病&未患病\\\midrule
有暴露&$a$&$b$\\无暴露&$c$&$d$\\\bottomrule\end{tabular}\end{center}
\subsection{什么是风险和比值？}
疾病风险为发病人数除以风险人数，即暴露者$a/(a+b)$，未暴露者$c/(c+d)$。比值为发病人数除以未发病人数，即$a/b$或$c/d$。
\subsection{比值比有两种算法，但只有一个比值比}
队列和临床试验的第一种算法（OR$_1$）为暴露者疾病比值除以未暴露者疾病比值：$\mathrm{OR}_1=(a/b)\div(c/d)=ad/bc$。病例对照的第二种算法（OR$_2$）为病例暴露比值除以对照暴露比值：$\mathrm{OR}_2=(a/c)\div(b/d)=ad/cb$。

因$ad/cb=ad/bc$，两种算法数值相同，只有一个比值比：$\mathrm{OR}_2=\mathrm{OR}_1$。
\subsection{疾病罕见时，比值近似风险，比值比近似风险比}
\sourcepage{274}
若两暴露组疾病都罕见，多数人不发病，风险人数近似未发病人数，所以罕见病的比值近似其风险。因此，比较两组疾病比值的OR$_1$也近似比较两组疾病风险的RR，即$\mathrm{OR}_1\approx\mathrm{RR}$。而病例对照可估计的病例暴露比值与对照暴露比值之比OR$_2$，也必须近似风险比：$\mathrm{OR}_2=\mathrm{OR}_1\approx\mathrm{RR}$。

\chapterappendix{9C}{第9章习题：病例对照研究的设计}
\sourcepage{275}
\begin{enumerate}
\item 第8章习题1c设计了维生素B12低与髋部骨折的回顾性队列。如何用同一队列更高效地开展病例对照，回答该问题？
\item 研究问题：“卵巢癌家族史使卵巢癌风险增加多少？”研究者计划病例对照。
\begin{enumerate}
\item 如何选择病例较好？
\item 据此如何选对照？
\item 病例与对照抽样可能有什么偏倚来源，方向如何？
\item 如何测量预测变量“卵巢癌家族史”？测量有何偏倚来源？
\item 用何种关联指标、何种统计学显著性检验？
\item 病例对照适合此问题吗？相对于其他可能设计，有何优缺点？
\end{enumerate}
\item 研究赛车类电子游戏与作为驾驶者发生真实车祸的风险关系。
\begin{enumerate}
\item 假定关注长期习惯性玩游戏的效应，病例对照如何选病例、对照，并测暴露？
\item 若关注驾驶前紧接一小时内玩游戏是否增加短期风险，应采用哪种间歇暴露短期效应设计？具体如何开展？
\end{enumerate}
\end{enumerate}

\endgroup
